% Generated draft: requires mathematical and source review.
\documentclass[11pt]{article}
\usepackage[margin=25mm]{geometry}
\usepackage{fontspec}
\setmainfont{Latin Modern Roman}
\usepackage{amsmath,amssymb,mathtools,unicode-math}
\setmathfont{Latin Modern Math}
\usepackage{enumitem,needspace,xurl,hyperref}
\setlength{\parindent}{0pt}
\setlength{\parskip}{5pt}
\setlength{\emergencystretch}{3em}
\newcommand{\N}{\mathbb{N}}
\newcommand{\Z}{\mathbb{Z}}
\newcommand{\Q}{\mathbb{Q}}
\newcommand{\R}{\mathbb{R}}
\newcommand{\C}{\mathbb{C}}
\newcommand{\F}{\mathbb{F}}
\newcommand{\A}{\mathbb{A}}
\newcommand{\PP}{\mathbb{P}}
\newcommand{\E}{\mathbb{E}}
\newcommand{\eps}{\varepsilon}
\newcommand{\dd}{\,\mathrm{d}}
\begin{document}
\section*{UnsolvedMath ID: 30006280}
\textbf{Source problem code:} catalogue entry
\subsection{Definitions and mathematical statement}
For a smooth projective variety \(W\), write \(D^b(W)=D^b(\operatorname{Coh}(W))\) for the bounded derived category of coherent sheaves. A Fano variety is a smooth projective variety whose anticanonical line bundle is ample. The source does not specify the ground field or its characteristic; these must be fixed before a fully specified conjecture is obtained. An extremal contraction is understood in the usual Mori-theoretic sense of a contraction associated with an extremal ray, subject to the unresolved qualifications concerning the omitted diagram below.

A semiorthogonal decomposition
\[
D^b(W)=\langle\mathcal C_1,\ldots,\mathcal C_m\rangle
\]
consists of full triangulated subcategories that generate \(D^b(W)\) and satisfy
\[
\operatorname{Hom}_{D^b(W)}(C_j,C_i[n])=0
\quad\text{for all }j>i,\ n\in\Z,\ C_j\in\mathcal C_j,\ C_i\in\mathcal C_i.
\]
A subcategory is admissible if its inclusion has both left and right adjoints. Mutations of admissible semiorthogonal blocks use these adjoints to replace adjacent blocks by the corresponding left or right mutated blocks. The admissibility hypotheses and the precise mutation requirement intended in the source are not explicitly supplied.

\textbf{Conjectural statement, with unresolved source qualifications.} Let \(X,Y,Z\) be smooth Fano varieties, and let
\[
f\colon X\longrightarrow Y,
\qquad h\colon X\longrightarrow Z
\]
be extremal contractions. The supplied question asks whether, for every such configuration in the intended class, there exist finite semiorthogonal decompositions
\[
D^b(Y)=\langle\mathcal A_1,\ldots,\mathcal A_s\rangle,
\qquad
D^b(Z)=\langle\mathcal B_1,\ldots,\mathcal B_t\rangle
\]
and additional blocks \(\mathcal P_1,\ldots,\mathcal P_r\) and \(\mathcal Q_1,\ldots,\mathcal Q_u\) in \(D^b(X)\) such that
\[
\begin{aligned}
D^b(X)
&=\langle\mathbf Lf^*\mathcal A_1,\ldots,\mathbf Lf^*\mathcal A_s,
\mathcal P_1,\ldots,\mathcal P_r\rangle,\\
D^b(X)
&=\langle\mathcal Q_1,\ldots,\mathcal Q_u,
\mathbf Lh^*\mathcal B_1,\ldots,\mathbf Lh^*\mathcal B_t\rangle,
\end{aligned}
\]
with the two decompositions of \(D^b(X)\) related by mutation. Here the displayed pullback blocks denote the essential images of the corresponding subcategories under derived pullback. This display makes explicit the natural block interpretation of compatibility suggested by the supplied excerpt; the excerpt itself suppresses the pullback notation and does not give a formal definition of compatibility. Whether the restrictions of pullback are required to identify the original blocks fully faithfully with admissible blocks, and whether the relation requires one mutation or a finite sequence of mutations, remain unspecified. The question is existential in the decompositions; the supplied wording does not assert compatibility for every preassigned decomposition of either target.

The original statement refers to contractions in an omitted diagram labelled \((1)\). Consequently, the supplied record does not establish whether that diagram imposes further restrictions on the contractions or varieties.

\textbf{Supplied motivating situation.} The excerpt also states
\[
X=M,\qquad Y=\mathbb P^{3g-3},\qquad Z=\operatorname{SU}_{C}(2,\Lambda),
\]
with \(Y=\mathbb P^{3g-2}\) and \(Z=\operatorname{SU}_{C}(2,\Lambda_0)\) in the even-degree case. It describes starting with a known decomposition of \(D^b(Y)\), extending it to \(D^b(X)\), and seeking mutations giving compatibility with a decomposition of \(D^b(Z)\). The objects \(M,C,\Lambda,\Lambda_0\), the parameter \(g\), and the relevant degree and smoothness assumptions are not defined in the supplied excerpt. These expressions therefore preserve the motivating notation without asserting a self-contained example satisfying all hypotheses.
\subsection{Short English statement}
The supplied conjecture asks whether two extremal contractions from a smooth Fano variety to smooth Fano targets admit semiorthogonal decompositions whose pullback blocks extend to decompositions of the source variety, with the resulting decompositions related by mutation. The ground field, possible restrictions in the omitted diagram, the exact meaning of compatibility, and the required mutation relation remain unresolved. The supplied open status is dated curation, not an independently verified conclusion.
\subsection{Sources}
\begin{itemize}
\item Local UnsolvedMath registry, permanent ID \(30006280\), source problem code \texttt{catalogue entry}, title \emph{Semiorthogonal Decompositions Along Extremal Contractions}; supplied complete source record associated with \texttt{data/unsolvedmath/problems.json} and the project registry \texttt{catalog/problems.json}. Attribution: Oberwolfach Reports Open Problems collection; proposed year \(2025\). No individual proposer is supplied.
\item Supplied original reference: \emph{Recent Trends in Algebraic Geometry} (2025), Oberwolfach Reports, DOI \texttt{10.4171/owr/2025/28}, pp.\ 1451--1516, \url{https://doi.org/10.4171/owr/2025/28}. The excerpt identifies the assertion as Conjecture 3. The report records workshop discussions and open problems; no narrower page location is supplied.
\item Supplied literature source, labelled ``original OWR report'': \url{https://doi.org/10.4171/owr/2025/28}. Literature assessment checked at \texttt{2026-08-22}: the record labels the conjecture open and states that known examples and mutations do not establish the general assertion. This is supplied dated curation derived from the OWR record, not an independently verified literature assessment; time-sensitive status and cited claims require independent verification.
\item The registry labels the statement \texttt{corrected\_verified} and reports that the official report was checked, the question isolated from surrounding exposition, and its notation repaired. These are supplied verification claims, not verification performed for this draft.
\end{itemize}
\end{document}
